\section{Conclusion}
\label{sec:conclusion}

Trained models of dynamics often reproduce hidden attractors but rarely their hiddenness. Degree theory fixes the bookkeeping
of a model's equilibria. When the true equilibria have an index sum different from one, a bounded model must either miss a
true equilibrium or add spurious ones; in an equilibrium-free system every bounded model has an equilibrium, and in all
our reproductions it made the training attractor self-excited. Otherwise
training decides how many unstable equilibria a model has, and an attractor is excited exactly through their unstable
manifolds. Across about $10^4$ models of four classes trained on six regimes, the probability that a reproduced hidden
attractor stays hidden is set mainly by the number of unstable equilibria in the region of the true system, with the class, the
reservoir size and the depth acting mainly through this number. Data from the basin of the visible attractor help; removing unstable equilibria during training keeps the attractor hidden but also removes true
ones; and generic models do not infer a hidden attractor that the data do not touch. Counting and testing the unstable
equilibria of a trained model is therefore a necessary step before its global predictions, and in particular its account of
hidden regimes, can be trusted.
